\documentclass[10pt]{article}
\usepackage[T1]{fontenc}
\usepackage[utf8]{inputenc}
\usepackage[french]{babel}
\usepackage{amsmath,amssymb,mathrsfs,mathpazo}
\usepackage{geometry}
\geometry{a4paper,margin=2cm}
\usepackage{graphicx,array,tabularx,tabularray,colortbl,multirow,multicol,booktabs,enumitem,pifont,ulem,fancybox,diagbox,cancel}
\usepackage{pstricks,pstricks-add,pst-plot,pst-tree,pst-node,pst-text,pst-func,pst-math,pst-bspline,pst-3dplot,pst-eucl,pst-solides3d}
\let\pscalculate\undefined
\usepackage{pst-bezier}
\usepackage{tikz,pgfplots,tkz-tab}
\usetikzlibrary{arrows,arrows.meta,patterns,calc,positioning,trees,decorations.markings,shapes.geometric,shapes.misc,shapes.arrows,intersections,decorations.pathmorphing,decorations.pathreplacing,backgrounds,shadows,babel}
\pgfplotsset{compat=1.18}
\usepgfplotslibrary{fillbetween}
\usepackage[np]{numprint}
\usepackage{listings,makecell,pas-tableur,twoopt}
\let\barre\undefined
\usepackage[pstricks]{tablvar}
\definecolor{darkgreen}{RGB}{0,100,0}
\providecommand{\red}{\color{red}}
\providecommand{\blue}{\color{blue}}
\providecommand{\green}{\color{green}}
\providecommand{\black}{\color{black}}
\providecommand{\white}{\color{white}}
\providecommand{\gray}{\color{gray}}
\providecommand{\cyan}{\color{cyan}}
\providecommand{\magenta}{\color{magenta}}
\providecommand{\yellow}{\color{yellow}}
\providecommand{\violet}{\color{violet}}
\providecommand{\orange}{\color{orange}}
\providecommand{\euro}{\texteuro}
\providecommand{\EUR}{\texteuro}
\providecommand{\ieme}{\textsuperscript{e}}
\providecommand{\degres}{\ensuremath{^\circ}}
\providecommand{\degree}{\ensuremath{^\circ}}
\providecommand{\vv}[1]{\overrightarrow{#1}}
\providecommand{\Coord}[1]{#1}
\providecommand{\decosix}{}
\providecommand{\starredbullet}{\ensuremath{\bullet}}
\pagestyle{empty}
\setlength{\parindent}{0pt}
\usepackage{adjustbox}\DeclareUnicodeCharacter{2236}{:}\DeclareUnicodeCharacter{2190}{\ensuremath{\leftarrow}}
\providecommand{\R}{\mathbb{R}}
\providecommand{\N}{\mathbb{N}}
\providecommand{\Z}{\mathbb{Z}}
\providecommand{\Q}{\mathbb{Q}}
\providecommand{\D}{\mathbb{D}}
\providecommand{\C}{\mathbb{C}}
\providecommand{\e}{{\mathrm{e}}}
\providecommand{\vect}[1]{\overrightarrow{#1}}
\providecommand{\vectt}[1]{\overrightarrow{\mathrm{#1}}}
\providecommand{\vvt}[1]{\overrightarrow{\mathrm{#1}}}
\providecommand{\np}[2][]{{#2~#1}}
\providecommand{\vv}[1]{{\overrightarrow{#1}}}
\renewcommand{\d}{{\mathrm{d}}}
\providecommand{\degres}{{^\circ}}
\providecommand{\degre}{{^\circ}}
\providecommand{\blue}{\color{blue}}
\providecommand{\red}{\color{red}}
\providecommand{\enskip}{\quad}
\providecommand{\llbracket}{\lbrack\!\lbrack}
\providecommand{\rrbracket}{\rbrack\!\rbrack}
\providecommand{\euro}{{\text{€}}}
\providecommand{\no}{{\text{n°}}}
\providecommand{\interval}[2]{{[#1;#2]}}
\providecommand{\intervFO}[2]{{[#1;#2[}}
\providecommand{\dotfill}{\ldots}

\providecommand{\ds}{\displaystyle}
\providecommand{\textding}[1]{#1}
\providecommand{\Oij}{\(\mathrm{O};\vec{\imath},\vec{\jmath}\)}
\providecommand{\Oijk}{\(\mathrm{O};\vec{\imath},\vec{\jmath},\vec{k}\)}
\providecommand{\Ouv}{\(\mathrm{O};\vec{u},\vec{v}\)}
\providecommand{\intervFF}[2]{\ensuremath{[#1;#2]}}
\renewcommand{\euro}{\ifmmode\text{\texteuro}\else\texteuro\fi}
\DeclareUnicodeCharacter{00B0}{\ifmmode^{\circ}\else\textdegree\fi}
\DeclareUnicodeCharacter{20AC}{\ifmmode\text{\texteuro}\else\texteuro\fi}
% Énoncés : Éducation nationale. Transcription et corrigés : APMEP, Denis Vergès et contributeurs. Découpage et classement : CoopMaths, Sébastien Lozano, Éric Elter et contributeurs. Adaptation de la présentation : Maths Explorer.
% CC BY-SA : https://coopmaths.fr/www/about/
% https://coopmaths.fr/alea/static/bac/2022/tex/bac_2022_05_sujet1_madagascar_4_cor.tex
\begin{document}
\noindent\textit{APMEP / CoopMaths. CC BY-SA. Présentation adaptée par Maths Explorer.}\par\medskip
\medskip

\begin{tabularx}{\linewidth}{|X|}\hline
\textbf{Principaux domaines abordés :} Géométrie dans l'espace\\ \hline
\end{tabularx}

\medskip

L'espace est muni d'un repère orthonormé \Oijk.

On considère les points A$(5~;~0~;~-1)$, B$(1~;~4~;~-1)$, C(1~;~0~;~3), D(5~;~4~;~3) et E(10~;~9~;~8).

\medskip

\begin{enumerate}
\item 
	\begin{enumerate}
		\item 
		

$\bullet~~$On a R$(3~;~2~;~-1)$ ;

$\bullet~~$$\vect{\text{AB}}\begin{pmatrix}-4\\4\\0\end{pmatrix}$.
		\item 
Si $\vect{\text{AB}}$ est un vecteur normal au plan $\mathcal{P}_1$ on sait qu'une équation de ce plan est :

$-4x + 4y + 0z = d$, avec $d \in \R$.

Comme R$(3~;~2~;~-1) \in \mathcal{P}_1 \iff -12 + 8 + 0 = d \iff d = - 4$.

Donc $M(x~;~y~;~z) \in \mathcal{P}_1 \iff - 4x + 4y = -4 \iff x - y - 1 = 0$.
		\item Démontrer que le point E appartient au plan $\mathcal{P}_1$ et que EA = EB.
On a E$(10~;~9~;~8) \in \mathcal{P}_1 \iff 10 - 9 - 1 = 0$ : cette égalité est vraie.

D'autre part $\vect{\text{EA}}\begin{pmatrix}-5\\-9\\-9\end{pmatrix}$ et 
$\vect{\text{EB}}\begin{pmatrix}-9\\-5\\-9\end{pmatrix}$.

On en déduit : EA$^2 = 25 + 81 + 81 = 187$ et EA$^2 = 81 + 25 + 81 = 187$.

EA$^2 = \text{EB}^2 \Rightarrow \text{EA} = \text{EB}$.

	\end{enumerate}	
\item On considère le plan $\mathcal{P}_2$ d'équation cartésienne $x - z - 2 = 0$.
	\begin{enumerate}
		\item 
On a vu que le plan $\mathcal{P}_1$ a pour vecteur normal $\vect{\text{AB}}\begin{pmatrix}-4\\4\\0\end{pmatrix}$.

Le plan $\mathcal{P}_2$ a pour vecteur normal $\vect{n}\begin{pmatrix}1\\0\\-1\end{pmatrix}$.

Ces deux vecteurs directeurs ne sont pas colinéaires, donc $\mathcal{P}_1$ et $\mathcal{P}_2$ ne sont pas parallèles, ils sont donc sécants.
		\item 
		






Si $M(x~;~y;~z)$ est commun aux deux plans ses coordonnées vérifient les équations des deux plans, donc le système :

$\left\{\begin{array}{l c l}
x - y - 1 &=& 0\\
x - z - 2 &=& 0
\end{array}\right.$. En posant $z = t$ le système devient :

$\left\{\begin{array}{l c l}
x - y - 1 &=& 0\\
x - z - 2 &=& 0\\
z&=&t
\end{array}\right.\Rightarrow \left\{\begin{array}{l c l}
x - 1 &=&y\\
x &=&z + 2\\
z&=&t
\end{array}\right. \iff \left\{\begin{array}{l c l}
x - 1 &=&y\\
x &=&t + 2\\
z&=&t
\end{array}\right. \iff$

$\left\{\begin{array}{l c l}
t + 2 - 1 &=&y\\
x &=&t + 2\\
z&=&t
\end{array}\right. 
\iff \left\{\begin{array}{l c l}
y&=&t + 1\\
x &=&t + 2\\
z&=&t
\end{array}\right.$, avec $t \in \R$.
	\end{enumerate}
\item  


Si $\Omega(x~;~y~;~z)$ est commun à la droite $\Delta$ et au plan $\mathcal{P}_3$ ses coordonnées vérifient l'équation du plan et les équations paramétriques de la droite, soit le système :

$\left\{\begin{array}{l c l}
y&=&t + 1\\
x &=&t + 2\\
z&=&t\\
y + z - 3&=&0
\end{array}\right.$. En remplaçant $x,\:y,\:z$ par leurs expressions en fonction de $t$ dans l'équation du plan on obtient :

$t + 1 + t - 3= 0 \iff 2t - 2 = 0 \iff 2t = 2 \iff t = 1$.

On a donc $\Omega\left(3~;~2~;~1 \right)$.
\end{enumerate}

\medskip

Si S et T sont deux points distincts de l'espace, on rappelle que l'ensemble des points M de l'espace tels que MS = MT est un plan, appelé plan médiateur du segment [ST]. On admet que les plans $\mathcal{P}_1$,\: $\mathcal{P}_2$ et $\mathcal{P}_3$ sont les plans médiateurs respectifs des segments [AB], [AC] et [AD].

\begin{enumerate}[resume]
\item
	\begin{enumerate}
		\item 
Sans admettre que les plans $\mathcal{P}_1$,\: $\mathcal{P}_2$ et $\mathcal{P}_3$ sont les plans médiateurs respectifs des segments [AB], [AC] et [AD], on peut calculer :

$\vect{\Omega \text{A}}\begin{pmatrix}2\\-2\\-2\end{pmatrix},\quad \vect{\Omega \text{B}}\begin{pmatrix}-2\\-2\\-2\end{pmatrix}, \quad \vect{\Omega \text{C}}\begin{pmatrix}-2\\-2\\2\end{pmatrix},\quad \vect{\Omega \text{D}}\begin{pmatrix}2\\2\\2\end{pmatrix}$.

On a donc $\Omega\text{A}^2 = 4 + 4 + 4 = 12, \: \Omega\text{B}^2 = 12,\: \Omega\text{C}^2 = 12, \: \Omega\text{D}^2 = 12$ et par conséquent :

$\Omega\text{A} = \Omega\text{B} = \Omega\text{C} = \Omega\text{D} = \sqrt{12} = \sqrt{4 \times 3} = 2\sqrt{3}$.
		\item 


Le résultat précédent montre que $\Omega$ est équidistant de A, B, C et D donc est le centre de la sphère de rayon $2\sqrt{3}$ contenant A, B, C et D.
	\end{enumerate}
\end{enumerate}

\end{document}
